% ==============================================================================
% CHAPTER 5: IMPLEMENTATION & EXPERIMENTAL SETUP
% Chair of Wireless Systems (Fachgebiet Drahtlose Systeme)
% ==============================================================================

\chapter{Implementation and Experimental Setup}
\label{ch:implementation}

\begin{takeawaybox}[Writing Guide: Chapter 5 --- Implementation]
\textbf{Goal}: Document the practical engineering realization of your system in sufficient detail to guarantee exact reproducibility by another researcher.\\
\textbf{Typical Length}: 10--15 pages (approx. 15--20\% of the thesis).\\
\textbf{Key Elements}:
\begin{itemize}[leftmargin=*]
    \item \textbf{Hardware Testbed}: Detailed specifications of edge boards, transceivers, processors, and sensors (\Cref{tab:hardware_specs}).
    \item \textbf{Software Stack}: Operating system, frameworks, libraries, quantization toolchains.
    \item \textbf{Dataset Details}: Source, train/val/test splits, sensor sampling frequencies, synthetic noise injection.
    \item \textbf{Code Snippets}: Concise, well-commented code listings illustrating non-trivial components (\Cref{lst:nli_clustering}).
\end{itemize}
\end{takeawaybox}

\section{Hardware and Testbed Configuration}
\label{sec:impl_hardware}

To validate \ac{UTSR} in realistic edge scenarios, experiments were conducted across a distributed laboratory testbed comprising two tiers of edge devices connected over IEEE 802.15.4 and Wi-Fi links. \Cref{tab:hardware_specs} summarizes the hardware specifications.

\begin{table}[ht]
    \centering
    \caption{Hardware specifications of the edge experimental testbed.}
    \label{tab:hardware_specs}
    \footnotesize
    \begin{tabularx}{\textwidth}{p{2.2cm}p{2.8cm}p{3.4cm}p{2.2cm}X}
        \toprule
        \textbf{Component} & \textbf{Device Model} & \textbf{Processor / Core} & \textbf{Memory} & \textbf{Role in Pipeline} \\
        \midrule
        Sensor Node & ESP32-WROOM-32 & Xtensa Dual-Core 32-bit & \SI{520}{\kibi\byte} SRAM & Ultrasonic telemetry capture \\
        Edge Gateway & NVIDIA Jetson Orin Nano & 6-core ARM Cortex-A78AE & \SI{8}{\gibi\byte} LPDDR5 & Local inference \& UQ routing \\
        Cloud Cluster & Server Node & Intel Xeon Gold 6330 & \SI{64}{\gibi\byte} RAM, A100 & High-tier escalation model \\
        \bottomrule
    \end{tabularx}
\end{table}

\section{Software Stack and Model Optimization}
\label{sec:impl_software}

The software pipeline was implemented in Python 3.11 utilizing PyTorch and HuggingFace Transformers. To meet edge latency requirements, the local language model was quantized to 4-bit precision via AWQ and compiled using ONNX Runtime. The bidirectional NLI scoring module relies on a distilled DeBERTa-v3-small model.

\Cref{lst:nli_clustering} illustrates the core Python implementation of the semantic equivalence graph construction.

\begin{lstlisting}[language=Python, caption={Core implementation of bidirectional semantic clustering.}, label={lst:nli_clustering}]
import torch
from transformers import AutoTokenizer, AutoModelForSequenceClassification

def build_semantic_adjacency(samples, nli_model, tokenizer, threshold=0.5):
    """
    Constructs a semantic adjacency matrix from generated sequences
    via bidirectional Natural Language Inference (NLI).
    """
    n = len(samples)
    adjacency = torch.eye(n, dtype=torch.bool)
    
    for i in range(n):
        for j in range(i + 1, n):
            # Check entailment in both directions: s_i -> s_j and s_j -> s_i
            pairs = [[samples[i], samples[j]], [samples[j], samples[i]]]
            inputs = tokenizer(pairs, padding=True, return_tensors="pt")
            with torch.no_grad():
                logits = nli_model(**inputs).logits
                probs = torch.softmax(logits, dim=-1)
                
            # Class index 2 corresponds to 'Entailment'
            fwd_entail = probs[0, 2].item() > threshold
            bwd_entail = probs[1, 2].item() > threshold
            
            if fwd_entail and bwd_entail:
                adjacency[i, j] = True
                adjacency[j, i] = True
                
    return adjacency
\end{lstlisting}

\section{Evaluation Datasets and Protocols}
\label{sec:impl_datasets}

We evaluate performance across three distinct benchmarks representing diverse uncertainty regimes:
\begin{enumerate}
    \item \textbf{HotpotQA}~\cite{yang2018hotpotqa}: A multi-hop reasoning benchmark where questions require chaining facts across Wikipedia articles. Characterized by predominantly \emph{epistemic} uncertainty.
    \item \textbf{AmbigQA}~\cite{min2020ambigqa}: An open-domain benchmark containing inherently ambiguous queries with multiple plausible answers. Characterized by high \emph{aleatoric} uncertainty.
    \item \textbf{Ultrasonic WSN Telemetry}: Real-world vibration and acoustic anomaly telemetry gathered from industrial sensor nodes at the Chair of Wireless Systems.
\end{enumerate}
